\usepackage[margin=1.05in]{geometry}
\usepackage{amsmath,amssymb,amsthm,mathtools}
\usepackage{array}
\usepackage{enumitem}
\usepackage{float}
\usepackage{hyperref}
\usepackage{xcolor}
\usepackage{tikz}
\usepackage{tikz-cd}

\usetikzlibrary{
  arrows.meta,
  calc,
  decorations.markings,
  decorations.pathmorphing,
  decorations.pathreplacing,
  positioning,
  shapes.geometric
}

\hypersetup{
  colorlinks=true,
  linkcolor=blue!55!black,
  urlcolor=blue!55!black,
  citecolor=blue!55!black
}

\makeatletter
\renewcommand{\@pnumwidth}{2.3em}
\renewcommand*\l@chapter[2]{%
  \ifnum \c@tocdepth >\m@ne
    \addpenalty{-\@highpenalty}%
    \vskip 1.0em \@plus\p@
    \setlength\@tempdima{2.8em}%
    \begingroup
      \parindent \z@ \rightskip \@pnumwidth
      \parfillskip -\@pnumwidth
      \leavevmode \bfseries
      \advance\leftskip\@tempdima
      \hskip -\leftskip
      #1\nobreak\hfil \nobreak\hb@xt@\@pnumwidth{\hss #2}\par
      \penalty\@highpenalty
    \endgroup
  \fi}
\makeatother

\setlength{\parindent}{0pt}
\setlength{\parskip}{0.55em}
\pagestyle{plain}
\setlist[itemize]{leftmargin=1.6em,itemsep=0.2em,topsep=0.2em}
\setlist[enumerate]{leftmargin=1.8em,itemsep=0.2em,topsep=0.2em}

\numberwithin{equation}{chapter}

\theoremstyle{definition}
\newtheorem{definition}{Definition}[chapter]
\newtheorem{assumption}[definition]{Assumption}
\newtheorem{construction}[definition]{Construction}
\newtheorem{convention}[definition]{Convention}
\newtheorem{example}[definition]{Example}
\newtheorem{hypothesis}[definition]{Hypothesis}
\newtheorem{controlledapproximation}[definition]{Controlled Approximation}

\theoremstyle{plain}
\newtheorem{lemma}[definition]{Lemma}
\newtheorem{proposition}[definition]{Proposition}
\newtheorem{theorem}[definition]{Theorem}
\newtheorem{quotedtheorem}[definition]{Quoted Theorem}
\newtheorem{corollary}[definition]{Corollary}
\newtheorem{conjecture}[definition]{Conjecture}

\theoremstyle{remark}
\newtheorem{remark}[definition]{Remark}
\newtheorem{warning}[definition]{Warning}
\newtheorem{openproblem}[definition]{Open Problem}

\newcommand{\dd}{\mathrm{d}}
\newcommand{\ii}{\mathrm{i}}
\newcommand{\ee}{\mathrm{e}}
\newcommand{\id}{\mathrm{id}}
\newcommand{\QFT}{\mathrm{QFT}}
\newcommand{\EFT}{\mathrm{EFT}}
\newcommand{\UV}{\mathrm{UV}}
\newcommand{\IR}{\mathrm{IR}}
\newcommand{\Hilb}{\mathcal H}
\newcommand{\Obs}{\mathcal A}
\newcommand{\Poinc}{\mathcal P_+^\uparrow}
\newcommand{\vac}{\Omega}
\newcommand{\ket}[1]{\lvert #1\rangle}
\newcommand{\bra}[1]{\langle #1\rvert}
\newcommand{\braket}[2]{\langle #1\mid #2\rangle}
\newcommand{\vev}[1]{\langle #1\rangle}
\newcommand{\R}{\mathbb R}
\newcommand{\C}{\mathbb C}

\newcommand{\ChapterConventions}[5]{%
  \par\smallskip
  \noindent\textbf{Conventions in this chapter.}
  \begin{itemize}[leftmargin=1.6em,itemsep=0pt,topsep=0.15em,parsep=0pt]
    \item \textbf{Signature.} #1.
    \item \textbf{Metric sign.} #2.
    \item \(\boldsymbol{\hbar}\). #3.
    \item \textbf{Vacuum symbol.} #4.
    \item \textbf{Generator convention.} #5.
  \end{itemize}
}
\newcommand{\ConventionsLorentzian}{%
  \ChapterConventions{Lorentzian \(\mathbb M^D\)}
  {\(\eta_{\mu\nu}=\operatorname{diag}(-,+,\ldots,+)\)}
  {\(1\), unless explicitly displayed}
  {\(\vac=\Omega\)}
  {\(U(a)=\exp(\ii a^\mu P_\mu)\), hence
  \([P_\mu,\widehat\Phi(x)]=\ii\partial_\mu\widehat\Phi(x)\) when the
  field is translation covariant}%
}
\newcommand{\ConventionsEuclidean}{%
  \ChapterConventions{Euclidean \(\R^D\)}
  {positive metric \(\delta_{\mu\nu}\)}
  {\(1\), unless explicitly displayed}
  {\(\vac=\Omega\) only after Hilbert-space reconstruction; otherwise
  brackets denote the stated functional expectation}
  {Lorentzian generators introduced by continuation use
  \(U(a)=\exp(\ii a^\mu P_\mu)\)}%
}
\newcommand{\ConventionsMixed}{%
  \ChapterConventions{Lorentzian formulas use \(\mathbb M^D\); Euclidean
  formulas use \(\R^D\)}
  {Lorentzian \(\eta_{\mu\nu}=\operatorname{diag}(-,+,\ldots,+)\);
  Euclidean \(\delta_{\mu\nu}\)}
  {\(1\), unless explicitly displayed}
  {\(\vac=\Omega\) for Hilbert-space vacuum matrix elements; functional
  averages are named separately}
  {Lorentzian translations use \(U(a)=\exp(\ii a^\mu P_\mu)\)}%
}
\newcommand{\ConventionsRealTime}{%
  \ChapterConventions{real time; no spacetime metric unless a field-theory
  example is explicitly introduced}
  {field-theory examples use
  \(\eta_{\mu\nu}=\operatorname{diag}(-,+,\ldots,+)\)}
  {\(1\), unless explicitly displayed}
  {\(\vac=\Omega\) when a vacuum vector is present}
  {time evolution is \(\exp(-\ii Ht)\), and spacetime translations use
  \(U(a)=\exp(\ii a^\mu P_\mu)\) when introduced}%
}
\newcommand{\ConventionsEuclideanCFT}{%
  \ChapterConventions{Euclidean \(\R^D\), with radial quantization invoked
  where stated}
  {positive metric \(\delta_{\mu\nu}\)}
  {\(1\), unless explicitly displayed}
  {\(\vac=\Omega\) for the radial-quantization vacuum when a Hilbert space is
  used}
  {represented symmetry transformations use \(\exp(\ii\epsilon^A\widehat Q_A)\);
  differential conformal generators are defined with the sign displayed in
  the relevant formula}%
}
\newcommand{\ConventionsMixedCFT}{%
  \ChapterConventions{Lorentzian local-operator data and Euclidean radial or
  conformal geometry where stated}
  {Lorentzian \(\eta_{\mu\nu}=\operatorname{diag}(-,+,\ldots,+)\);
  Euclidean \(\delta_{\mu\nu}\)}
  {\(1\), unless explicitly displayed}
  {\(\vac=\Omega\) for the Lorentzian or radial-quantization vacuum}
  {Poincare translations use \(U(a)=\exp(\ii a^\mu P_\mu)\); represented
  conformal charges use \(\exp(\ii\epsilon^A\widehat Q_A)\) when a unitary
  Hilbert-space action is present}%
}
